\documentclass[../DoAn.tex]{subfiles}
\begin{document}

Chương này trình bày kết quả khảo sát hiện trạng công tác thi đua -- khen thưởng tại Học viện Khoa học Quân sự (HVKHQS) và rút ra các yêu cầu chức năng, phi chức năng của phần mềm. Mục \ref{section:2.1} mô tả thực trạng và các nhóm khó khăn. Mục \ref{section:2.2} trình bày cấu trúc tác nhân, sơ đồ use case tổng quát và tám sơ đồ phân rã. Mục \ref{section:2.3} đặc tả sáu use case trọng yếu. Mục \ref{section:2.4} tổng hợp các yêu cầu phi chức năng về hiệu năng, bảo mật và tính nhất quán dữ liệu.

\section{Khảo sát hiện trạng}
\label{section:2.1}

Trước khi có phần mềm chuyên dụng, công tác thi đua -- khen thưởng tại HVKHQS dựa vào ba nguồn dữ liệu: hồ sơ giấy lưu tại Phòng Chính trị, các tệp Microsoft Excel chia sẻ qua thư mục mạng nội bộ và sổ theo dõi tay của từng đơn vị. Qua trao đổi với cán bộ phụ trách thi đua, nhóm thực hiện đồ án ghi nhận một số đặc điểm sau.

Cấp đề xuất đầu tiên là cán bộ chỉ huy đơn vị hoặc cán bộ chính trị cấp cơ quan, nắm danh sách quân nhân thuộc phạm vi quản lý và lập đề nghị theo từng năm hoặc theo các đợt đặc biệt. Cấp xét duyệt là Phòng Chính trị Học viện, kiểm tra điều kiện rồi lập tờ trình gửi Ban Giám đốc. Cấp ký quyết định là Giám đốc Học viện với các danh hiệu thuộc thẩm quyền, hoặc cấp trên (Tổng cục II, Bộ Quốc phòng, Thủ tướng Chính phủ) với các danh hiệu vượt thẩm quyền. Quân nhân được khen thưởng tiếp nhận thông tin thụ động qua văn bản hoặc bảng tin đơn vị.

Một quân nhân thường có nhiều bản ghi thuộc các nhóm khác nhau: chuỗi danh hiệu hằng năm tích luỹ trong toàn bộ thời gian phục vụ, Huy chương Chiến sĩ vẻ vang (HCCSVV) ở mốc 10/15/20 năm, Huân chương Bảo vệ Tổ quốc (HCBVTQ), Kỷ niệm chương, Huy chương Quân kỳ quyết thắng và các khen thưởng đột xuất. Cán bộ phục vụ trên 25 năm có thể có 30--40 bản ghi; nhân với quy mô vài nghìn cán bộ, học viên và sĩ quan, dữ liệu lịch sử ước lượng vài trăm nghìn bản ghi.

Bốn vấn đề được phản ánh nhiều nhất qua khảo sát:

\begin{enumerate}
    \item Tra cứu lịch sử khen thưởng cá nhân chậm do phải mở và đối chiếu nhiều tệp Excel.
    \item Đối chiếu chuỗi danh hiệu hằng năm khó kiểm soát do khung xét lùi theo từng năm: chu kỳ có thể lặp lại, BKTTCP cá nhân chỉ được nhận một lần, BKBTBQP của chu kỳ trước rơi ngoài khung 3 năm khi xét CSTĐTQ chu kỳ mới.
    \item Không có cơ chế phát hiện trùng đề xuất khi nhiều cán bộ cùng phụ trách một đơn vị.
    \item Nhật ký kiểm toán gần như không tồn tại, gây khó truy hồi khi phát hiện sai sót.
\end{enumerate}

Các đặc thù nghiệp vụ trên đòi hỏi một hệ thống có thể (a) lưu trữ tập trung và liên kết các loại bản ghi, (b) tự động kiểm tra điều kiện theo bộ quy tắc đã mã hoá, (c) phân quyền theo cây tổ chức và (d) ghi nhật ký kiểm toán cho mọi thao tác làm thay đổi dữ liệu nghiệp vụ.

\section{Tổng quan chức năng}
\label{section:2.2}

Phần này trình bày cấu trúc tác nhân và các nhóm chức năng chính của hệ thống thông qua sơ đồ use case tổng quát và các sơ đồ phân rã chi tiết.

\subsection{Sơ đồ use case tổng quát}
\label{subsection:2.2.1}

Hệ thống có bốn nhóm tác nhân tương ứng với bốn cấp phân quyền nội bộ:

\begin{itemize}
    \item \textbf{SuperAdmin} đại diện cho bộ phận kỹ thuật và công nghệ thông tin. Nhóm này không tham gia trực tiếp vào nghiệp vụ xét khen thưởng mà tập trung vào quản trị hạ tầng, sao lưu -- khôi phục dữ liệu, theo dõi nhật ký an ninh và quản trị các tài khoản cấp quản trị.
    \item \textbf{Admin} là cán bộ thuộc Phòng Chính trị Học viện. Nhóm này có toàn quyền nghiệp vụ: tạo và phê duyệt đề xuất, quản lý quân nhân, đơn vị, chức vụ, quản lý kho file PDF quyết định, gắn số quyết định cho đề xuất khi phê duyệt và xuất các báo cáo.
    \item \textbf{Manager} là cán bộ chỉ huy hoặc cán bộ chính trị cấp đơn vị, phụ trách lập đề xuất khen thưởng cho phạm vi đơn vị mình quản lý và theo dõi trạng thái phê duyệt. Manager có thể thao tác trên dữ liệu quân nhân thuộc đơn vị nhưng không thể chuyển quân nhân giữa các đơn vị (chức năng này chỉ thuộc Admin).
    \item \textbf{User} là cán bộ -- học viên -- sĩ quan đăng nhập để tra cứu hồ sơ khen thưởng cá nhân, xem các mốc danh hiệu hằng năm đang đến kỳ đề nghị và nhận thông báo về kết quả phê duyệt.
\end{itemize}

Hệ thống cung cấp 15 nhóm chức năng chính: Đăng nhập, Quản lý tài khoản, Quản lý quân nhân, Quản lý đơn vị, Khen thưởng hằng năm, Khen thưởng HCCSVV, Khen thưởng HCBVTQ, Khen thưởng thành tích khoa học, Khen thưởng đột xuất, Đề xuất khen thưởng, Kiểm tra điều kiện và gợi ý, Thông báo thời gian thực, Xem nhật ký hệ thống, Sao lưu và khôi phục, Báo cáo và thống kê. Mối quan hệ giữa từng tác nhân và các nhóm chức năng được thể hiện ở Hình \ref{fig:usecase-tongquat}.

\begin{figure}[H]
    \centering
    % \includegraphics[width=0.9\textwidth]{Hinhve/Hinh-2-1-use-case-tong-quat.png}
    \caption{Sơ đồ use case tổng quát hệ thống PM QLKT}
    \label{fig:usecase-tongquat}
\end{figure}

\subsection{Phân rã: Quản lý quân nhân và lịch sử chức vụ}
\label{subsection:2.2.2}

Quân nhân là thực thể trung tâm, liên kết tới hầu hết các bảng nghiệp vụ (danh hiệu hằng năm, HCCSVV, HCBVTQ, thành tích khoa học, khen thưởng đột xuất, đề xuất). Module này có tần suất tương tác cao nhất hệ thống. Phân rã chức năng gồm các use case con:

\begin{itemize}
    \item \textbf{Tạo mới quân nhân}: chỉ Admin có quyền. Yêu cầu nhập tối thiểu họ tên, số CCCD, ngày sinh, giới tính, ngày nhập ngũ, đơn vị (CQĐV hoặc ĐVTT), chức vụ. Hệ thống kiểm tra trùng CCCD trước khi ghi vào cơ sở dữ liệu.
    \item \textbf{Cập nhật thông tin}: Admin sửa được toàn bộ; Manager chỉ sửa được các trường ngoại trừ đơn vị và CCCD; bản thân quân nhân (User) chỉ sửa được địa chỉ, số điện thoại, email.
    \item \textbf{Chuyển đơn vị}: chỉ Admin. Khi chuyển, hệ thống cập nhật số lượng quân nhân của đơn vị cũ (giảm 1) và đơn vị mới (tăng 1) trong cùng một transaction. Quan trọng là toàn bộ lịch sử khen thưởng được giữ nguyên, không tạo bản ghi mới hay xoá bản ghi cũ.
    \item \textbf{Quản lý lịch sử chức vụ}: Admin nhập danh sách các giai đoạn quân nhân giữ chức vụ với hệ số tương ứng. Dữ liệu này là đầu vào cho thuật toán xét khen thưởng HCBVTQ.
    \item \textbf{Nhập / xuất Excel}: cho phép tải lên hoặc tải xuống danh sách quân nhân theo định dạng quy chuẩn. Thao tác nhập đi qua hai bước (xem trước, xác nhận) sẽ được mô tả chi tiết tại \ref{subsection:2.2.7}.
\end{itemize}

\begin{figure}[H]
    \centering
    \caption{Sơ đồ use case phân rã chức năng Quản lý quân nhân}
    \label{fig:usecase-quannhan}
\end{figure}

\subsection{Phân rã: Quản lý đơn vị và chức vụ}
\label{subsection:2.2.3}

Cấu trúc tổ chức của Học viện được mô hình hoá thành cây hai cấp. Bảng \texttt{CoQuanDonVi} (CQĐV) lưu các đơn vị cấp cha; bảng \texttt{DonViTrucThuoc} (ĐVTT) lưu các đơn vị con thuộc một CQĐV duy nhất. Khi nhập thông tin quân nhân, một trong hai trường \texttt{co\_quan\_don\_vi\_id} hoặc \texttt{don\_vi\_truc\_thuoc\_id} sẽ được điền, không bao giờ điền cả hai cùng lúc.

Module quản lý đơn vị bao gồm các use case con: tạo mới CQĐV, tạo mới ĐVTT thuộc một CQĐV cụ thể, cập nhật thông tin, xoá (chỉ thực hiện được khi không còn quân nhân nào thuộc đơn vị) và xem cây tổ chức dưới dạng tree-view trên giao diện. Bảng \texttt{ChucVu} quản lý danh mục chức vụ kèm hệ số \texttt{he\_so\_chuc\_vu} (giá trị từ 0.5 đến 1.2 tuỳ chức danh) và liên kết tới CQĐV hoặc ĐVTT mà chức vụ đó thuộc về. Hệ số chức vụ là dữ liệu đầu vào trực tiếp cho thuật toán tính tháng cống hiến theo nhóm hệ số 0.7 / 0.8 / 0.9--1.0.

\begin{figure}[H]
    \centering
    \caption{Sơ đồ use case phân rã chức năng Quản lý đơn vị và chức vụ}
    \label{fig:usecase-donvi}
\end{figure}

\subsection{Phân rã: Đề xuất khen thưởng theo bảy nhóm}
\label{subsection:2.2.4}

Đề xuất khen thưởng là điểm khởi đầu của mọi quy trình ghi nhận thành tích. Hệ thống hỗ trợ bảy loại đề xuất riêng biệt, mỗi loại có dữ liệu đầu vào và quy tắc xét điều kiện khác nhau:

\begin{enumerate}
    \item \textbf{Cá nhân hằng năm} (\texttt{CA\_NHAN\_HANG\_NAM}): danh hiệu CSTT, CSTĐCS, BKBTBQP, CSTĐTQ, BKTTCP cộng các cờ tương ứng và số quyết định cấp cao hơn.
    \item \textbf{Đơn vị hằng năm} (\texttt{DON\_VI\_HANG\_NAM}): danh hiệu ĐVQT, ĐVTT cộng các cờ BKBTBQP / BKTTCP cấp đơn vị.
    \item \textbf{HCCSVV} (\texttt{NIEN\_HAN}): Huy chương Chiến sĩ vẻ vang theo các mốc 10 / 15 / 20 năm phục vụ tích luỹ tới năm xét.
    \item \textbf{HCBVTQ} (\texttt{CONG\_HIEN}): Huân chương Bảo vệ Tổ quốc dựa trên thời gian giữ chức vụ thuộc các nhóm hệ số khác nhau.
    \item \textbf{Kỷ niệm chương} (\texttt{KNC\_VSNXD\_QDNDVN}): Kỷ niệm chương vì sự nghiệp xây dựng QĐNDVN, một lần duy nhất theo giới tính (nam 25 năm phục vụ, nữ 20 năm).
    \item \textbf{Quân kỳ quyết thắng} (\texttt{HC\_QKQT}): Huy chương Quân kỳ quyết thắng cho cán bộ phục vụ đủ 25 năm.
    \item \textbf{Thành tích khoa học} (\texttt{NCKH}): Khen thưởng đề tài khoa học hoặc sáng kiến kinh nghiệm.
\end{enumerate}

Cả bảy loại dùng chung tầng dữ liệu là bảng \texttt{BangDeXuat} với các trường JSON \texttt{data\_danh\_hieu}, \texttt{data\_thanh\_tich}, \texttt{data\_nien\_han}, \texttt{data\_cong\_hien} để lưu chi tiết riêng cho từng loại. Lựa chọn JSON cho phép schema linh hoạt theo loại đề xuất mà không phải bổ sung bảng trung gian. Mỗi loại được gắn với một thực thể tuân theo giao diện \texttt{ProposalStrategy} đảm nhiệm xây dựng payload, kiểm tra và nhập dữ liệu khi phê duyệt; chi tiết trình bày tại Chương \ref{chapter:Experiment}.

Trên giao diện, quy trình tạo đề xuất gồm ba bước: chọn loại đề xuất và năm/tháng $\rightarrow$ chọn quân nhân hoặc đơn vị $\rightarrow$ thiết lập danh hiệu cho từng đối tượng. Bước hai cho phép chọn hàng loạt qua các bộ lọc theo đơn vị, năm sinh, giới tính và mốc thời gian phục vụ.

\begin{figure}[H]
    \centering
    \caption{Sơ đồ use case phân rã chức năng Đề xuất khen thưởng}
    \label{fig:usecase-dexuat}
\end{figure}

\subsection{Phân rã: Danh hiệu hằng năm cá nhân}
\label{subsection:2.2.5}

Đây là module nghiệp vụ phức tạp nhất hệ thống: quy tắc danh hiệu hằng năm đi kèm khung xét lùi theo từng năm, lại có một danh hiệu chỉ được nhận đúng một lần trong toàn bộ thời gian phục vụ. Danh hiệu hằng năm cá nhân gồm ba cấp BKBTBQP -- CSTĐTQ -- BKTTCP với các điều kiện sau:

\begin{itemize}
    \item \textbf{BKBTBQP}: chu kỳ 2 năm, yêu cầu 2 năm CSTĐCS liên tục và NCKH mỗi năm; không có danh hiệu tiền điều kiện; có thể nhận lặp lại sau mỗi chu kỳ.
    \item \textbf{CSTĐTQ}: chu kỳ 3 năm, yêu cầu đúng 1 BKBTBQP trong 3 năm gần nhất tính từ năm trước năm xét và NCKH mỗi năm; có thể nhận lặp lại sau mỗi chu kỳ.
    \item \textbf{BKTTCP}: chu kỳ 7 năm, yêu cầu 7 năm CSTĐCS liên tục (độ dài chuỗi là bội số của 7), đúng 3 BKBTBQP và đúng 2 CSTĐTQ trong 7 năm gần nhất, NCKH mỗi năm; \textbf{chỉ nhận một lần duy nhất} với cá nhân.
\end{itemize}

Use case chính của module là \textbf{Kiểm tra điều kiện và gợi ý}. Hệ thống tính từ dữ liệu lịch sử rồi sinh ra hai loại thông tin cho mỗi quân nhân: (a) danh sách danh hiệu đủ điều kiện đề nghị trong năm hiện tại kèm lý do bằng tiếng Việt; (b) các chỉ số khung xét -- số BKBTBQP trong 3 năm hoặc 7 năm gần nhất, số CSTĐTQ trong 7 năm gần nhất, độ dài chuỗi CSTĐCS đang đếm. Cán bộ Phòng Chính trị tham chiếu các gợi ý này để chốt danh sách đề xuất.

Use case phụ là \textbf{Tính lại tổng thể}: khi một bản ghi danh hiệu hằng năm cũ thay đổi, hệ thống chạy lại thuật toán cho riêng quân nhân đó hoặc cho toàn bộ, giữ cho bảng suy diễn \texttt{HoSoHangNam} luôn nhất quán với bảng nguồn \texttt{DanhHieuHangNam}.

\begin{figure}[H]
    \centering
    \caption{Sơ đồ use case phân rã chức năng khen thưởng cá nhân hằng năm}
    \label{fig:usecase-chuoi-canhan}
\end{figure}

\subsection{Phân rã: Danh hiệu hằng năm đơn vị}
\label{subsection:2.2.6}

Danh hiệu hằng năm đơn vị có cấu trúc đơn giản hơn cá nhân do không có cấp CSTĐTQ và không yêu cầu thành tích NCKH. Danh hiệu hằng năm đơn vị chỉ gồm hai cấp:

\begin{itemize}
    \item \textbf{BKBTBQP đơn vị} chu kỳ 2 năm, yêu cầu 2 năm ĐVQT liên tục, không có danh hiệu tiền điều kiện, có thể nhận lặp lại sau mỗi chu kỳ.
    \item \textbf{BKTTCP đơn vị} chu kỳ 7 năm, yêu cầu ít nhất 3 BKBTBQP đơn vị trong 7 năm gần nhất, có thể nhận lặp lại sau mỗi chu kỳ.
\end{itemize}

Khác biệt quan trọng với danh hiệu hằng năm cá nhân là điều kiện ``ít nhất 3 BKBTBQP'' thay vì ``đúng 3 BKBTBQP''. Vì đơn vị có thể nhận lặp lại BKTTCP sau mỗi chu kỳ 7 năm, ngưỡng đếm cần linh hoạt hơn so với danh hiệu một lần duy nhất của cá nhân.

\begin{figure}[H]
    \centering
    \caption{Sơ đồ use case phân rã chức năng khen thưởng đơn vị hằng năm}
    \label{fig:usecase-chuoi-donvi}
\end{figure}

\subsection{Phân rã: Nhập và xuất dữ liệu Excel}
\label{subsection:2.2.7}

Lượng dữ liệu lịch sử của Học viện hiện lưu trên Excel; chức năng nhập Excel hàng loạt là cầu nối để hệ thống tiếp nhận dữ liệu cũ thay vì nhập tay từng bản ghi. Quy trình gồm bốn use case con:

\begin{itemize}
    \item \textbf{Tải tệp mẫu}: hệ thống sinh tệp Excel có cấu trúc cột chuẩn theo từng loại nghiệp vụ (quân nhân, danh hiệu hằng năm, lịch sử chức vụ, HCCSVV, ...) cùng các Data Validation rule cho phép chọn từ danh sách thả xuống ngay trong tệp.
    \item \textbf{Xem trước}: sau khi cán bộ điền dữ liệu và tải tệp lên, hệ thống đọc nội dung từng dòng, đối chiếu với schema Zod và các quy tắc nghiệp vụ, sau đó hiển thị kết quả gồm danh sách dòng hợp lệ và danh sách dòng có lỗi (kèm chỉ số dòng và mô tả lỗi). Bước này không ghi vào cơ sở dữ liệu.
    \item \textbf{Xác nhận}: cán bộ kiểm tra danh sách hợp lệ rồi xác nhận. Hệ thống mở một transaction Prisma, ghi tuần tự từng dòng; nếu xảy ra lỗi tại bất kỳ dòng nào, toàn bộ thao tác được hủy bỏ và cơ sở dữ liệu giữ nguyên trạng thái trước khi nhập.
    \item \textbf{Xuất danh sách}: cho phép xuất kết quả khen thưởng theo nhiều tiêu chí (theo năm, theo đơn vị, theo loại danh hiệu) ra tệp Excel có định dạng phù hợp với mẫu báo cáo nội bộ của Học viện.
\end{itemize}

\begin{figure}[H]
    \centering
    \caption{Sơ đồ use case phân rã chức năng Nhập -- xuất Excel}
    \label{fig:usecase-excel}
\end{figure}

\subsection{Phân rã: Quản trị hệ thống}
\label{subsection:2.2.8}

Module quản trị bao quát các chức năng nền tảng phục vụ vận hành hệ thống. Bao gồm:

\begin{itemize}
    \item \textbf{Quản lý tài khoản}: SuperAdmin tạo tài khoản Admin; Admin tạo tài khoản Manager và User; mỗi tài khoản gắn với một quân nhân duy nhất qua khoá ngoại \texttt{quan\_nhan\_id}.
    \item \textbf{Đặt lại mật khẩu}: Admin có thể đặt lại mật khẩu cho Manager / User về giá trị mặc định.
    \item \textbf{Nhật ký hệ thống} (\texttt{SystemLog}): ghi lại mọi thao tác thay đổi dữ liệu kèm thông tin gồm thời điểm, người thực hiện, loại tài nguyên, mã định danh tài nguyên, mô tả thao tác và dữ liệu trước, sau khi thay đổi.
    \item \textbf{Sao lưu và khôi phục}: SuperAdmin cấu hình lịch sao lưu định kỳ qua khu vực DevZone; hệ thống dùng \texttt{pg\_dump} sinh tệp SQL và lưu tại thư mục \texttt{backups/}. Thao tác khôi phục cũng được thực hiện qua giao diện này.
    \item \textbf{Cấu hình hệ thống} (\texttt{SystemSetting}): bật/tắt các cờ tính năng, ví dụ \texttt{allow\_notify\_import} cho phép gửi thông báo khi nhập Excel hàng loạt.
\end{itemize}

Tất cả các thao tác trong module quản trị đều có lưu nhật ký riêng để đảm bảo trách nhiệm cá nhân và truy hồi.

\begin{figure}[H]
    \centering
    \caption{Sơ đồ use case phân rã chức năng Quản trị hệ thống}
    \label{fig:usecase-quantri}
\end{figure}

\subsection{Quy trình nghiệp vụ --- ba luồng quan trọng}
\label{subsection:2.2.9}

Bên cạnh sơ đồ use case, ba quy trình nghiệp vụ phức tạp nhất được mô hình hoá thêm bằng sơ đồ hoạt động (activity diagram) để làm rõ luồng xử lý từng bước.

\textbf{Quy trình A: Đề xuất $\rightarrow$ Phê duyệt một đợt khen thưởng.} Manager chọn loại đề xuất và năm xét, hệ thống lọc danh sách quân nhân/đơn vị thuộc phạm vi quản lý. Manager chọn các đối tượng đủ điều kiện, thiết lập danh hiệu cụ thể rồi gửi đề xuất. Hệ thống lưu vào bảng \texttt{BangDeXuat} ở trạng thái \texttt{PENDING} và phát thông báo thời gian thực tới Admin. Admin mở danh sách chờ duyệt, đối chiếu từng đối tượng và bổ sung số quyết định trước khi phê duyệt; nếu cần thay đổi hạng danh hiệu, Admin từ chối kèm lý do để Manager đề xuất lại với hạng đúng. Khi phê duyệt, hệ thống mở một transaction để (a) chuyển trạng thái đề xuất sang \texttt{APPROVED}, (b) ghi các bản ghi danh hiệu/khen thưởng vào bảng tương ứng (\texttt{DanhHieuHangNam}, \texttt{KhenThuongHCCSVV}, \dots), (c) gắn số quyết định và đường dẫn tệp PDF, (d) ghi nhật ký kiểm toán, (e) thông báo tới Manager đã đề xuất và (f) tính lại các hồ sơ suy diễn liên quan.

\textbf{Quy trình B: Tính điều kiện khen thưởng hằng năm cho một quân nhân.} Khi cán bộ truy cập trang hồ sơ của một quân nhân, hệ thống nạp toàn bộ bản ghi \texttt{DanhHieuHangNam} của quân nhân đó, sắp xếp giảm dần theo năm và chạy hàm \texttt{computeChainContext}. Hàm này tính ba chỉ số: (a) độ dài chuỗi CSTĐCS hiện tại (\texttt{streakCstdcs}), (b) năm gần nhất có cờ BKBTBQP / CSTĐTQ / BKTTCP, (c) số cờ tiền điều kiện nằm trong khung xét tương ứng (3 năm gần nhất cho CSTĐTQ, 7 năm gần nhất cho BKTTCP). Sau đó với mỗi cấp trong \texttt{PERSONAL\_CHAIN\_AWARDS}, hàm \texttt{checkChainEligibility} đối chiếu số liệu thực tế với cấu hình cấp và trả về \texttt{eligible: boolean} kèm thông điệp \texttt{reason} bằng tiếng Việt.

\textbf{Quy trình C: Nhập tệp Excel hàng loạt.} Cán bộ Admin tải lên tệp Excel chứa danh sách bản ghi cần nhập. Hệ thống đọc tệp ở backend, tách thành các sheet (vd: \texttt{QuanNhan}, \texttt{DanhHieuHangNam}, \texttt{ThanhTichKhoaHoc}) và đọc từng dòng. Mỗi dòng được kiểm tra qua schema Zod và các quy tắc nghiệp vụ liên quan (vd: CCCD không được trùng, đơn vị phải tồn tại). Các dòng hợp lệ và lỗi được tổng hợp vào kết quả \texttt{ImportPreview} trả về frontend. Cán bộ xem trước, đối chiếu, có thể quay lại sửa tệp Excel rồi tải lên lại. Khi xác nhận, backend mở transaction Prisma, ghi tuần tự các dòng hợp lệ; nếu có bất kỳ lỗi nào tại thời điểm ghi, transaction được rollback hoàn toàn.

\section{Đặc tả use case chi tiết}
\label{section:2.3}

Phần này đặc tả sáu use case trọng yếu của hệ thống. Mỗi use case được trình bày qua hai bảng: bảng đặc tả luồng sự kiện (tên, tác nhân, tiền điều kiện, luồng chính, luồng thay thế, hậu điều kiện) và bảng dữ liệu đầu vào (các trường người dùng nhập). Sáu use case được lựa chọn vì phản ánh các nghiệp vụ phức tạp nhất và có nhiều tương tác giữa các module. Đặc tả các use case còn lại được trình bày tại Phụ lục \ref{appendix:usecase}.

\begin{table}[H]
\centering
\caption{Đặc tả use case UC-01 Đăng nhập}
\label{tab:uc01}
\renewcommand{\arraystretch}{1.2}
\begin{tabular}{|>{\centering\arraybackslash}p{3.2cm}|>{\centering\arraybackslash}p{0.8cm}|p{2.8cm}|p{6.4cm}|}
\hline
\textbf{Tên use case} & \multicolumn{3}{l|}{Đăng nhập hệ thống} \\ \hline
\textbf{Tác nhân} & \multicolumn{3}{l|}{SuperAdmin, Admin, Manager, User} \\ \hline
\textbf{Tiền điều kiện} & \multicolumn{3}{p{10.6cm}|}{Người dùng đã có tài khoản hợp lệ; trình duyệt cho phép \texttt{localStorage}.} \\ \hline
\multirow{7}{*}{\parbox{3.2cm}{\centering\textbf{Luồng sự kiện chính}\\\textbf{(Thành công)}}}
 & \textbf{STT} & \textbf{Thực hiện bởi} & \multicolumn{1}{c|}{\textbf{Hành động}} \\ \cline{2-4}
 & 1. & Người dùng & Truy cập trang \texttt{/login} \\ \cline{2-4}
 & 2. & Người dùng & Nhập tên đăng nhập và mật khẩu \\ \cline{2-4}
 & 3. & Người dùng & Bấm nút ``Đăng nhập'' \\ \cline{2-4}
 & 4. & Hệ thống & Đối chiếu mật khẩu băm bằng bcrypt \\ \cline{2-4}
 & 5. & Hệ thống & Phát hành Access Token và Refresh Token \\ \cline{2-4}
 & 6. & Hệ thống & Chuyển hướng tới trang chủ tương ứng vai trò \\ \hline
\multirow{4}{*}{\parbox{3cm}{\centering\textbf{Luồng sự kiện}\\\textbf{thay thế}}}
 & \textbf{STT} & \textbf{Thực hiện bởi} & \multicolumn{1}{c|}{\textbf{Hành động}} \\ \cline{2-4}
 & 4a. & Hệ thống & Báo ``Tài khoản hoặc mật khẩu không đúng'' \\ \cline{2-4}
 & 4b. & Hệ thống & Báo ``Tài khoản đã bị tạm khoá, liên hệ quản trị viên'' \\ \cline{2-4}
 & 5a. & Hệ thống & Vượt ngưỡng 30 yêu cầu đăng nhập thất bại / IP / 5 phút, từ chối tạm thời với mã 429 \\ \hline
\textbf{Hậu điều kiện} & \multicolumn{3}{p{10.6cm}|}{Phiên đăng nhập được thiết lập; nhật ký đăng nhập được ghi với \texttt{action = LOGIN}.} \\ \hline
\end{tabular}
\end{table}

\begin{table}[H]
\centering
\caption{Dữ liệu đầu vào cho chức năng đăng nhập}
\label{tab:uc01-input}
\begin{tabular}{|c|p{3.5cm}|p{4.5cm}|c|p{2.5cm}|}
\hline
\textbf{STT} & \textbf{Trường dữ liệu} & \textbf{Mô tả} & \textbf{Bắt buộc} & \textbf{Ví dụ hợp lệ} \\ \hline
1 & Tên đăng nhập & Tên tài khoản đăng nhập & Có & \texttt{ducta} \\ \hline
2 & Mật khẩu & Mật khẩu đăng nhập & Có & \texttt{Dat\#2026} \\ \hline
\end{tabular}
\end{table}

\begin{table}[H]
\centering
\caption{Đặc tả use case UC-02 Tạo mới quân nhân}
\label{tab:uc02}
\renewcommand{\arraystretch}{1.2}
\begin{tabular}{|>{\centering\arraybackslash}p{3.2cm}|>{\centering\arraybackslash}p{0.8cm}|p{2.8cm}|p{6.4cm}|}
\hline
\textbf{Tên use case} & \multicolumn{3}{l|}{Tạo mới hồ sơ quân nhân} \\ \hline
\textbf{Tác nhân} & \multicolumn{3}{l|}{Admin} \\ \hline
\textbf{Tiền điều kiện} & \multicolumn{3}{p{10.6cm}|}{Đã đăng nhập với vai trò Admin; cây tổ chức (CQĐV/ĐVTT) đã có dữ liệu.} \\ \hline
\multirow{8}{*}{\parbox{3.2cm}{\centering\textbf{Luồng sự kiện chính}\\\textbf{(Thành công)}}}
 & \textbf{STT} & \textbf{Thực hiện bởi} & \multicolumn{1}{c|}{\textbf{Hành động}} \\ \cline{2-4}
 & 1. & Admin & Truy cập ``Quản lý quân nhân'' và bấm ``Thêm mới'' \\ \cline{2-4}
 & 2. & Admin & Điền các trường bắt buộc trong biểu mẫu \\ \cline{2-4}
 & 3. & Admin & Bấm ``Lưu'' \\ \cline{2-4}
 & 4. & Hệ thống & Xác thực dữ liệu hai phía bằng schema thống nhất \\ \cline{2-4}
 & 5. & Hệ thống & Kiểm tra trùng CCCD trong cơ sở dữ liệu \\ \cline{2-4}
 & 6. & Hệ thống & Ghi bản ghi mới và tăng số lượng quân nhân của đơn vị \\ \cline{2-4}
 & 7. & Hệ thống & Trả về kết quả thành công, cập nhật bảng danh sách \\ \hline
\multirow{4}{*}{\parbox{3cm}{\centering\textbf{Luồng sự kiện}\\\textbf{thay thế}}}
 & \textbf{STT} & \textbf{Thực hiện bởi} & \multicolumn{1}{c|}{\textbf{Hành động}} \\ \cline{2-4}
 & 4a. & Hệ thống & Có trường bắt buộc bị thiếu, chặn ngay tại biểu mẫu \\ \cline{2-4}
 & 5a. & Hệ thống & CCCD đã tồn tại, trả lỗi ``Số CCCD đã tồn tại trong hệ thống'' \\ \cline{2-4}
 & 6a. & Hệ thống & Đơn vị được chọn không tồn tại, trả lỗi và huỷ bỏ giao dịch \\ \hline
\textbf{Hậu điều kiện} & \multicolumn{3}{p{10.6cm}|}{Bản ghi mới có trong bảng quân nhân; số lượng quân nhân của đơn vị tăng 1; nhật ký \texttt{action = CREATE, resource = personnel} được ghi đầy đủ.} \\ \hline
\end{tabular}
\end{table}

\begin{table}[H]
\centering
\caption{Dữ liệu đầu vào cho chức năng tạo mới quân nhân}
\label{tab:uc02-input}
\begin{tabular}{|c|p{3.5cm}|p{4.5cm}|c|p{2.5cm}|}
\hline
\textbf{STT} & \textbf{Trường dữ liệu} & \textbf{Mô tả} & \textbf{Bắt buộc} & \textbf{Ví dụ hợp lệ} \\ \hline
1 & Họ và tên & Họ tên đầy đủ của quân nhân & Có & Trần Anh Đức \\ \hline
2 & CCCD & Số căn cước công dân (12 chữ số) & Có & 001202012345 \\ \hline
3 & Ngày sinh & Ngày tháng năm sinh & Có & 12/03/2002 \\ \hline
4 & Giới tính & Nam hoặc Nữ & Có & Nam \\ \hline
5 & Ngày nhập ngũ & Ngày bắt đầu phục vụ & Có & 01/09/2020 \\ \hline
6 & Đơn vị & CQĐV hoặc ĐVTT mà quân nhân thuộc về & Có & Khoa CNTT \\ \hline
7 & Chức vụ & Chức vụ hiện tại & Có & Trợ lý \\ \hline
\end{tabular}
\end{table}

\begin{table}[H]
\centering
\caption{Đặc tả use case UC-03 Tạo đề xuất khen thưởng}
\label{tab:uc03}
\renewcommand{\arraystretch}{1.2}
\begin{tabular}{|>{\centering\arraybackslash}p{3.2cm}|>{\centering\arraybackslash}p{0.8cm}|p{2.8cm}|p{6.4cm}|}
\hline
\textbf{Tên use case} & \multicolumn{3}{l|}{Tạo đề xuất khen thưởng cho một đợt} \\ \hline
\textbf{Tác nhân} & \multicolumn{3}{l|}{Admin, Manager} \\ \hline
\textbf{Tiền điều kiện} & \multicolumn{3}{p{10.6cm}|}{Đã đăng nhập; có ít nhất một quân nhân thuộc phạm vi quản lý của Manager.} \\ \hline
\multirow{9}{*}{\parbox{3.2cm}{\centering\textbf{Luồng sự kiện chính}\\\textbf{(Thành công)}}}
 & \textbf{STT} & \textbf{Thực hiện bởi} & \multicolumn{1}{c|}{\textbf{Hành động}} \\ \cline{2-4}
 & 1. & Người dùng & Truy cập ``Đề xuất khen thưởng'' và bấm ``Tạo đề xuất mới'' \\ \cline{2-4}
 & 2. & Người dùng & Bước 1: chọn loại đề xuất và năm/tháng xét \\ \cline{2-4}
 & 3. & Hệ thống & Lọc và hiển thị danh sách quân nhân/đơn vị đủ điều kiện \\ \cline{2-4}
 & 4. & Người dùng & Bước 2: tích chọn các đối tượng đề xuất \\ \cline{2-4}
 & 5. & Người dùng & Bước 3: thiết lập danh hiệu cụ thể cho từng đối tượng \\ \cline{2-4}
 & 6. & Người dùng & Tải lên các tệp đính kèm (nếu có) và bấm ``Gửi đề xuất'' \\ \cline{2-4}
 & 7. & Hệ thống & Xác thực dữ liệu, tạo bản ghi đề xuất ở trạng thái chờ duyệt \\ \cline{2-4}
 & 8. & Hệ thống & Phát thông báo thời gian thực tới các Admin liên quan \\ \hline
\multirow{4}{*}{\parbox{3cm}{\centering\textbf{Luồng sự kiện}\\\textbf{thay thế}}}
 & \textbf{STT} & \textbf{Thực hiện bởi} & \multicolumn{1}{c|}{\textbf{Hành động}} \\ \cline{2-4}
 & 4a. & Hệ thống & Có quân nhân chưa đủ điều kiện, hiển thị cảnh báo nhưng vẫn cho phép gửi nếu Manager xác nhận \\ \cline{2-4}
 & 4b. & Hệ thống & Manager chọn quân nhân ngoài phạm vi đơn vị, từ chối với lỗi 403 \\ \cline{2-4}
 & 7a. & Hệ thống & Có lỗi mạng, lưu nháp ở \texttt{localStorage} để không mất dữ liệu \\ \hline
\textbf{Hậu điều kiện} & \multicolumn{3}{p{10.6cm}|}{Bản ghi đề xuất được tạo ở trạng thái chờ duyệt; thông báo đã gửi tới Admin; nhật ký \texttt{action = CREATE, resource = proposals} được ghi cùng tóm tắt loại và số lượng đối tượng.} \\ \hline
\end{tabular}
\end{table}

\begin{table}[H]
\centering
\caption{Dữ liệu đầu vào cho chức năng tạo đề xuất khen thưởng}
\label{tab:uc03-input}
\begin{tabular}{|c|p{3.5cm}|p{4.5cm}|c|p{2.5cm}|}
\hline
\textbf{STT} & \textbf{Trường dữ liệu} & \textbf{Mô tả} & \textbf{Bắt buộc} & \textbf{Ví dụ hợp lệ} \\ \hline
1 & Loại đề xuất & Một trong bảy nhóm khen thưởng & Có & CSTĐCS \\ \hline
2 & Năm xét & Năm xét khen thưởng & Có & 2026 \\ \hline
3 & Tháng xét & Tháng xét (đối với đợt giữa năm) & Không & 12 \\ \hline
4 & Danh sách đối tượng & Quân nhân hoặc đơn vị được đề xuất & Có & 25 quân nhân \\ \hline
5 & Danh hiệu cụ thể & Hạng/cấp danh hiệu cho từng đối tượng & Có & Hạng Ba \\ \hline
6 & Tệp đính kèm & Tệp PDF biên bản, ảnh giấy khen & Không & \texttt{bb-2026.pdf} \\ \hline
\end{tabular}
\end{table}

\begin{table}[H]
\centering
\caption{Đặc tả use case UC-04 Phê duyệt đề xuất}
\label{tab:uc04}
\renewcommand{\arraystretch}{1.2}
\begin{tabular}{|>{\centering\arraybackslash}p{3.2cm}|>{\centering\arraybackslash}p{0.8cm}|p{2.8cm}|p{6.4cm}|}
\hline
\textbf{Tên use case} & \multicolumn{3}{l|}{Phê duyệt một đề xuất khen thưởng} \\ \hline
\textbf{Tác nhân} & \multicolumn{3}{l|}{Admin} \\ \hline
\textbf{Tiền điều kiện} & \multicolumn{3}{p{10.6cm}|}{Đề xuất đang ở trạng thái chờ duyệt; Admin có quyền phê duyệt loại đề xuất tương ứng.} \\ \hline
\multirow{8}{*}{\parbox{3.2cm}{\centering\textbf{Luồng sự kiện chính}\\\textbf{(Thành công)}}}
 & \textbf{STT} & \textbf{Thực hiện bởi} & \multicolumn{1}{c|}{\textbf{Hành động}} \\ \cline{2-4}
 & 1. & Admin & Mở chi tiết đề xuất từ danh sách ``Chờ duyệt'' \\ \cline{2-4}
 & 2. & Admin & Đối chiếu danh sách đối tượng và danh hiệu \\ \cline{2-4}
 & 3. & Admin & Nhập số quyết định, tháng -- năm nhận; có thể đính kèm tệp PDF quyết định \\ \cline{2-4}
 & 4. & Admin & Bấm ``Phê duyệt'' \\ \cline{2-4}
 & 5. & Hệ thống & Mở giao dịch: chuyển trạng thái, ghi danh hiệu, gắn quyết định, tính lại hồ sơ suy diễn \\ \cline{2-4}
 & 6. & Hệ thống & Phát thông báo thời gian thực tới Manager đã đề xuất và các quân nhân được khen \\ \cline{2-4}
 & 7. & Hệ thống & Ghi nhật ký kèm dữ liệu trước -- sau khi thay đổi \\ \hline
\multirow{5}{*}{\parbox{3cm}{\centering\textbf{Luồng sự kiện}\\\textbf{thay thế}}}
 & \textbf{STT} & \textbf{Thực hiện bởi} & \multicolumn{1}{c|}{\textbf{Hành động}} \\ \cline{2-4}
 & 3a. & Admin & Phát hiện hạng danh hiệu sai $\rightarrow$ chuyển sang luồng từ chối kèm lý do (5b); Manager đề xuất lại với hạng đúng \\ \cline{2-4}
 & 4a. & Hệ thống & Số quyết định trùng, cảnh báo và yêu cầu xác nhận trước khi tiếp tục \\ \cline{2-4}
 & 5b. & Admin & Từ chối thay vì phê duyệt: trạng thái chuyển ``Từ chối'' kèm lý do \\ \cline{2-4}
 & 5c. & Hệ thống & Lỗi tại bước ghi (vd: vi phạm quy tắc một lần duy nhất), huỷ toàn bộ giao dịch \\ \hline
\textbf{Hậu điều kiện} & \multicolumn{3}{p{10.6cm}|}{Trạng thái đề xuất là ``Đã duyệt'' hoặc ``Từ chối''; nếu được duyệt, các bản ghi khen thưởng đã có trong cơ sở dữ liệu, hồ sơ suy diễn được tính lại, thông báo đã gửi, nhật ký đầy đủ.} \\ \hline
\end{tabular}
\end{table}

\begin{table}[H]
\centering
\caption{Dữ liệu đầu vào cho chức năng phê duyệt đề xuất}
\label{tab:uc04-input}
\begin{tabular}{|c|p{3.5cm}|p{4.5cm}|c|p{2.5cm}|}
\hline
\textbf{STT} & \textbf{Trường dữ liệu} & \textbf{Mô tả} & \textbf{Bắt buộc} & \textbf{Ví dụ hợp lệ} \\ \hline
1 & Số quyết định & Số quyết định khen thưởng (tham chiếu PDF đã tải lên trước) & Có & 1234/QĐ-HV \\ \hline
2 & Tháng -- năm nhận & Tháng và năm gắn với quyết định khen thưởng & Có & 12/2025 \\ \hline
3 & Lý do từ chối & Áp dụng khi từ chối đề xuất & Khi từ chối & Hạng đề xuất chưa phù hợp \\ \hline
\end{tabular}
\end{table}

\begin{table}[H]
\centering
\caption{Đặc tả use case UC-05 Tính lại điều kiện khen thưởng và sinh gợi ý}
\label{tab:uc05}
\renewcommand{\arraystretch}{1.2}
\begin{tabular}{|>{\centering\arraybackslash}p{3.2cm}|>{\centering\arraybackslash}p{0.8cm}|p{2.8cm}|p{6.4cm}|}
\hline
\textbf{Tên use case} & \multicolumn{3}{l|}{Tính lại điều kiện khen thưởng hằng năm và sinh gợi ý} \\ \hline
\textbf{Tác nhân} & \multicolumn{3}{p{10.6cm}|}{Admin, Manager (kích hoạt chủ động); Hệ thống (tự động sau mỗi lần phê duyệt).} \\ \hline
\textbf{Tiền điều kiện} & \multicolumn{3}{p{10.6cm}|}{Quân nhân tồn tại; có ít nhất một bản ghi danh hiệu hằng năm.} \\ \hline
\multirow{6}{*}{\parbox{3.2cm}{\centering\textbf{Luồng sự kiện chính}\\\textbf{(Thành công)}}}
 & \textbf{STT} & \textbf{Thực hiện bởi} & \multicolumn{1}{c|}{\textbf{Hành động}} \\ \cline{2-4}
 & 1. & Hệ thống & Truy vấn toàn bộ danh hiệu hằng năm của quân nhân \\ \cline{2-4}
 & 2. & Hệ thống & Dẫn xuất ngữ cảnh danh hiệu hằng năm (số năm liên tục, danh hiệu trong khung xét) \\ \cline{2-4}
 & 3. & Hệ thống & Với mỗi cấp danh hiệu, gọi hàm xét điều kiện \\ \cline{2-4}
 & 4. & Hệ thống & Ghi kết quả vào bảng hồ sơ suy diễn \\ \cline{2-4}
 & 5. & Hệ thống & Trả thông điệp gợi ý bằng tiếng Việt cho frontend hiển thị \\ \hline
\multirow{4}{*}{\parbox{3cm}{\centering\textbf{Luồng sự kiện}\\\textbf{thay thế}}}
 & \textbf{STT} & \textbf{Thực hiện bởi} & \multicolumn{1}{c|}{\textbf{Hành động}} \\ \cline{2-4}
 & 3a. & Hệ thống & Quân nhân đã nhận BKTTCP, trả về thông điệp ``Đã có BKTTCP'' \\ \cline{2-4}
 & 3b. & Hệ thống & Số năm CSTĐCS liên tục chưa đủ chu kỳ, gợi ý số năm cần tích luỹ thêm \\ \cline{2-4}
 & 3c. & Hệ thống & Năm hiện tại đúng vào mốc lỡ đợt, giải thích cách xét lại ở chu kỳ kế tiếp \\ \hline
\textbf{Hậu điều kiện} & \multicolumn{3}{p{10.6cm}|}{Hồ sơ suy diễn nhất quán với dữ liệu nguồn; trang chi tiết quân nhân hiển thị các danh hiệu đủ điều kiện kèm thông điệp gợi ý.} \\ \hline
\end{tabular}
\end{table}

\begin{table}[H]
\centering
\caption{Đặc tả use case UC-06 Nhập dữ liệu khen thưởng hàng loạt từ tệp Excel}
\label{tab:uc06}
\renewcommand{\arraystretch}{1.2}
\begin{tabular}{|>{\centering\arraybackslash}p{3.2cm}|>{\centering\arraybackslash}p{0.8cm}|p{2.8cm}|p{6.4cm}|}
\hline
\textbf{Tên use case} & \multicolumn{3}{l|}{Nhập dữ liệu khen thưởng hàng loạt từ tệp Excel} \\ \hline
\textbf{Tác nhân} & \multicolumn{3}{l|}{Admin} \\ \hline
\textbf{Tiền điều kiện} & \multicolumn{3}{p{10.6cm}|}{Đã có tệp mẫu tải xuống từ hệ thống và điền dữ liệu theo cấu trúc cột quy định.} \\ \hline
\multirow{9}{*}{\parbox{3.2cm}{\centering\textbf{Luồng sự kiện chính}\\\textbf{(Thành công)}}}
 & \textbf{STT} & \textbf{Thực hiện bởi} & \multicolumn{1}{c|}{\textbf{Hành động}} \\ \cline{2-4}
 & 1. & Admin & Truy cập chức năng nhập tệp Excel và chọn loại nghiệp vụ \\ \cline{2-4}
 & 2. & Admin & Tải lên tệp \texttt{.xlsx} qua biểu mẫu \\ \cline{2-4}
 & 3. & Hệ thống & Đọc tệp, kiểm tra cấu trúc cột và đọc tuần tự từng dòng \\ \cline{2-4}
 & 4. & Hệ thống & Kiểm tra quy tắc nghiệp vụ và tìm quân nhân tương ứng cho mỗi dòng \\ \cline{2-4}
 & 5. & Hệ thống & Trả về bảng xem trước với các dòng hợp lệ và dòng lỗi kèm mô tả \\ \cline{2-4}
 & 6. & Admin & Kiểm tra danh sách lỗi, nếu đồng ý bấm ``Xác nhận nhập'' \\ \cline{2-4}
 & 7. & Hệ thống & Mở giao dịch và ghi tuần tự các dòng hợp lệ vào cơ sở dữ liệu \\ \cline{2-4}
 & 8. & Hệ thống & Hoàn tất giao dịch và thông báo kết quả thời gian thực \\ \hline
\multirow{4}{*}{\parbox{3cm}{\centering\textbf{Luồng sự kiện}\\\textbf{thay thế}}}
 & \textbf{STT} & \textbf{Thực hiện bởi} & \multicolumn{1}{c|}{\textbf{Hành động}} \\ \cline{2-4}
 & 3a. & Hệ thống & Tệp sai định dạng hoặc thiếu trang tính bắt buộc, báo lỗi và dừng xử lý \\ \cline{2-4}
 & 4a. & Admin & Có dòng lỗi ở bước xem trước, có thể tải tệp xuống sửa và tải lại \\ \cline{2-4}
 & 7a. & Hệ thống & Có lỗi tại bước ghi, huỷ bỏ toàn bộ giao dịch và giữ nguyên dữ liệu cũ \\ \hline
\textbf{Hậu điều kiện} & \multicolumn{3}{p{10.6cm}|}{Các bản ghi hợp lệ trong tệp được ghi vào cơ sở dữ liệu; nhật ký hành động được ghi với tổng kết số dòng thành công và số dòng lỗi; Admin nhận thông báo kết quả thời gian thực.} \\ \hline
\end{tabular}
\end{table}

\begin{table}[H]
\centering
\caption{Dữ liệu đầu vào cho chức năng nhập Excel hàng loạt}
\label{tab:uc06-input}
\begin{tabular}{|c|p{3.5cm}|p{4.5cm}|c|p{2.5cm}|}
\hline
\textbf{STT} & \textbf{Trường dữ liệu} & \textbf{Mô tả} & \textbf{Bắt buộc} & \textbf{Ví dụ hợp lệ} \\ \hline
1 & Loại nghiệp vụ & Quân nhân, danh hiệu hằng năm, lịch sử chức vụ, ... & Có & Danh hiệu \\ \hline
2 & Tệp Excel & Tệp \texttt{.xlsx} điền theo mẫu & Có & \texttt{cstdcs-2025.xlsx} \\ \hline
\end{tabular}
\end{table}

\section{Yêu cầu phi chức năng}
\label{section:2.4}

Bên cạnh các yêu cầu chức năng đã trình bày ở các mục trước, hệ thống cần đáp ứng một số yêu cầu phi chức năng để đảm bảo vận hành ổn định, an toàn và phù hợp với đặc thù môi trường quân sự. Các yêu cầu này được nhóm lại thành ba khía cạnh: hiệu năng và trải nghiệm người dùng, bảo mật, tính nhất quán dữ liệu.

\subsection{Hiệu năng và trải nghiệm người dùng}
\label{subsection:2.4.1}

Hệ thống cần vận hành ổn định trong mạng nội bộ Học viện với cao điểm vài chục người dùng đồng thời (đợt xét khen thưởng cuối năm). Các tiêu chí định lượng đặt ra như sau.

Thời gian phản hồi của các trang chính (trang chủ, danh sách quân nhân, danh sách đề xuất) phải dưới hai giây trên kết nối LAN bình thường. Tính lại hồ sơ khen thưởng hằng năm cho một quân nhân phải dưới một giây; tính lại tổng thể cho 3.000 quân nhân được phép kéo dài đến vài phút nhưng cần kèm thanh tiến trình phát qua Socket.IO.

Các bảng dữ liệu lớn (toàn bộ quân nhân, danh sách khen thưởng theo năm) phân trang phía máy chủ, mặc định 20 dòng/trang và tối đa 100 dòng/trang. Với các thao tác trạng thái thường gặp (đánh dấu đã đọc thông báo, tích chọn quân nhân vào đề xuất), giao diện cập nhật lạc quan ngay trước khi máy chủ phản hồi để giảm cảm giác chờ.

Dữ liệu được truy vấn thường xuyên nhưng ít thay đổi -- cây CQĐV/ĐVTT, danh mục chức vụ, danh mục danh hiệu -- được nạp một lần vào React Context và chia sẻ cho toàn ứng dụng, tránh trùng lặp yêu cầu giữa các thành phần.

\subsection{Bảo mật}
\label{subsection:2.4.2}

Hệ thống lưu dữ liệu cán bộ và sĩ quan trong môi trường quân sự nên bảo mật là yêu cầu ưu tiên hàng đầu. Bảy nhóm biện pháp được áp dụng song song.

\textbf{Thứ nhất: Xác thực bằng JWT có cơ chế làm mới luân phiên.} Mỗi phiên đăng nhập sinh một cặp Access Token (30 phút) và Refresh Token (7 ngày). Frontend lưu hai token vào \texttt{localStorage} và đính Access Token vào header \texttt{Authorization: Bearer} cho mỗi yêu cầu. Backend lưu thêm Refresh Token vào cột \texttt{TaiKhoan.refreshToken} để giữ khả năng thu hồi từ phía máy chủ. Middleware \texttt{verifyToken} đối chiếu cả chữ ký JWT lẫn giá trị trong cơ sở dữ liệu; xoá cột này tương đương buộc đăng xuất. Khi Access Token hết hạn, frontend gọi \texttt{/auth/refresh}; backend phát hành cặp token mới và ghi đè Refresh Token cũ, vô hiệu hoá mọi lần dùng lại token cũ. Cơ chế hạn chế thiệt hại nếu một Refresh Token bị đánh cắp.

\textbf{Thứ hai: Băm mật khẩu bằng bcrypt với hệ số chi phí 10.} Mật khẩu không lưu ở dạng văn bản thô. Hệ số 10 cân bằng giữa thời gian xác thực (khoảng 80\,ms trên máy chủ phát triển) và độ khó tấn công vét cạn nếu cơ sở dữ liệu bị rò rỉ \cite{bcrypt}.

\textbf{Thứ ba: Phân quyền bốn cấp qua middleware \texttt{checkRole}.} Mỗi route nhạy cảm đi qua middleware kiểm tra vai trò người dùng nằm trong danh sách cho phép. Bên cạnh kiểm tra vai trò, middleware \texttt{unitFilter} giới hạn phạm vi truy vấn của Manager về cây đơn vị mình phụ trách.

\textbf{Thứ tư: Xác thực dữ liệu hai phía bằng Zod.} Frontend dùng Zod \cite{zoddocs} xác thực ngay tại form để chặn dữ liệu sai định dạng từ trước khi gửi đi. Backend dùng cùng thư viện Zod xác thực lại tại điểm cuối qua middleware \texttt{validate(schema)} -- nếu yêu cầu đi thẳng tới API mà bỏ qua frontend, backend vẫn từ chối dữ liệu không hợp lệ. Phương thức \texttt{.strict()} loại bỏ các trường ngoài schema, ngăn người dùng đẩy lén thuộc tính không khai báo. Dùng chung thư viện ở hai phía giảm chi phí học, giảm trùng lặp khái niệm và tạo điều kiện trích xuất schema dùng chung trong tương lai.

\textbf{Thứ năm: Bảo vệ tầng vận chuyển và tầng ứng dụng.} Khi triển khai sản xuất, toàn bộ kết nối qua HTTPS (TLS 1.2 trở lên). Header bảo mật do middleware Helmet áp đặt (Content-Security-Policy, X-Frame-Options, X-Content-Type-Options). CORS được cấu hình chỉ chấp nhận tên miền frontend chính. Rate limiter chia hai cấp: \texttt{authLimiter} giới hạn 30 yêu cầu / IP / 5 phút cho các điểm cuối xác thực (đăng nhập, làm mới token, đặt lại mật khẩu) với cờ \texttt{skipSuccessfulRequests} để không chặn nhầm người dùng hợp lệ; \texttt{writeLimiter} giới hạn 30 yêu cầu / IP / 15 phút cho các thao tác ghi nhạy cảm.

\textbf{Thứ sáu: Nhật ký kiểm toán đầy đủ.} Mọi thao tác làm thay đổi dữ liệu nghiệp vụ đều ghi vào bảng \texttt{SystemLog} qua middleware \texttt{auditLog}. Tập hành động chuẩn hoá trong \texttt{AUDIT\_ACTIONS} gồm 17 mã: \texttt{CREATE}, \texttt{UPDATE}, \texttt{DELETE}, \texttt{IMPORT}, \texttt{IMPORT\_PREVIEW}, \texttt{EXPORT}, \texttt{APPROVE}, \texttt{REJECT}, \texttt{PROPOSE}, \texttt{RECALCULATE}, \texttt{BULK}, \texttt{BULK\_BYPASS}, \texttt{BACKUP}, \texttt{LOGIN}, \texttt{LOGOUT}, \texttt{CHANGE\_PASSWORD}, \texttt{RESET\_PASSWORD}. Trong đó \texttt{BULK\_BYPASS} đánh dấu các thao tác sửa dữ liệu cũ do SuperAdmin thực hiện qua tuyến \texttt{POST /api/awards/bulk-bypass} bỏ qua kiểm tra điều kiện -- mọi lần gọi đều được ghi log riêng phục vụ kiểm toán hậu kiểm. Mỗi bản ghi mang theo thời điểm, mã người thực hiện, vai trò, loại tài nguyên (\texttt{resource}), mã định danh tài nguyên (\texttt{tai\_nguyen\_id}) và mô tả tiếng Việt do hàm builder sinh tự động. Nhật ký thuộc \texttt{resource = backup} chỉ SuperAdmin truy cập được; các vai trò thấp hơn bị filter ngay tại tầng service. Tuyến \texttt{DELETE /api/system-logs} cũng giới hạn cho riêng SuperAdmin, đảm bảo Admin không thể xoá vết kiểm toán các thao tác của chính mình.

\textbf{Thứ bảy: Sao lưu định kỳ và khôi phục.} Tác vụ sao lưu chạy bên trong tiến trình backend qua \texttt{node-cron}: ở mỗi lần kích hoạt, dịch vụ đọc 21 bảng dữ liệu, sinh chuỗi câu lệnh \texttt{INSERT INTO ... VALUES (\dots)} đã escape, ghi tệp \texttt{.sql} vào thư mục \texttt{backups/} với quy ước tên \texttt{YYYY-MM-DD\_HH-mm-ss.sql}. Lịch chạy mặc định là \texttt{0 1 1 * *} (01:00 ngày 1 hằng tháng) và có thể chỉnh trực tiếp trong DevZone. Tệp cũ hơn \texttt{backup\_retention\_days} (mặc định 15) bị xoá ở bước cleanup. SuperAdmin tải xuống, xoá hoặc khôi phục thủ công qua giao diện DevZone (khôi phục thực thi bằng lệnh \texttt{psql} ngoài hệ thống).

\subsection{Tính nhất quán và chính xác của dữ liệu}
\label{subsection:2.4.3}

Một sai lệch nhỏ giữa các bản ghi có thể dẫn tới quyết định sai (vd: bỏ sót quân nhân đủ điều kiện, đề xuất nhầm chu kỳ); vì vậy tính nhất quán dữ liệu được đặt ngang hàng với bảo mật.

\textbf{Giao dịch ACID cho nghiệp vụ trên nhiều thực thể.} Các nghiệp vụ có thao tác ghi/đọc trên nhiều thực thể (bảng) cùng lúc được đóng gói trong transaction Prisma. Khi phê duyệt một đề xuất gồm 50 quân nhân, hệ thống thực hiện trong một transaction: 50 lần ghi \texttt{DanhHieuHangNam}, cập nhật trạng thái đề xuất, gắn số quyết định, ghi nhật ký, tính lại hồ sơ suy diễn. Hoặc tất cả thành công, hoặc cơ sở dữ liệu trở về trạng thái trước phê duyệt nếu một bước gặp lỗi.

\textbf{Kiểm tra điều kiện chuỗi ở hai tầng.} Hàm \texttt{checkChainEligibility} được gọi cả trong hàm tính lại hồ sơ (\texttt{computeEligibilityFlags}) lẫn trong hàm kiểm tra phê duyệt (\texttt{checkAwardEligibility}). Hai tầng phải đồng nhất: nếu xảy ra sai lệch, một phía chấp nhận đề xuất trong khi phía kia hiển thị ``không đủ điều kiện'', gây ảnh hưởng đến tính nhất quán của dữ liệu và mức độ tin cậy của hệ thống đối với người dùng. Cả hai hàm cùng tham chiếu một bộ kịch bản kiểm chứng thủ công để giữ đồng bộ.

\textbf{Tính lại hồ sơ suy diễn sau mỗi thay đổi nguồn.} Các bảng \texttt{HoSoHangNam}, \texttt{HoSoNienHan}, \texttt{HoSoCongHien} lưu kết quả tính toán (cờ đủ điều kiện, thông điệp gợi ý). Mọi thao tác tạo, sửa hoặc xoá bản ghi nguồn đều kích hoạt hàm tính lại tương ứng bên trong cùng transaction, đảm bảo dữ liệu suy diễn không bao giờ ``đi sau'' dữ liệu nguồn.

\textbf{Khoá ngoại nhất quán qua chính sách Prisma.} Quan hệ giữa các bảng khai báo qua \texttt{@relation} kèm chính sách \texttt{onDelete: Restrict} ở các quan hệ trọng yếu (vd: không cho xoá \texttt{QuanNhan} khi vẫn còn bản ghi \texttt{DanhHieuHangNam}). Một số bảng phụ trợ như thông báo dùng \texttt{onDelete: Cascade} để dọn dẹp tự động khi xoá đối tượng cha.

\textbf{Ràng buộc duy nhất qua chỉ mục unique.} Các trường yêu cầu duy nhất -- \texttt{cccd} của quân nhân, \texttt{username} của tài khoản, tổ hợp (\texttt{quan\_nhan\_id}, \texttt{nam}) của \texttt{DanhHieuHangNam} -- khai báo \texttt{@unique} ở schema. Khi vi phạm, Prisma trả lỗi \texttt{P2002}; backend bắt và chuyển thành thông điệp tiếng Việt thân thiện (vd: ``Quân nhân này đã có bản ghi danh hiệu năm 2024, vui lòng chỉnh sửa thay vì tạo mới'').


\end{document}
